\section{Elicitação e Estruturação de Variáveis mediante Método Delphi}

Este módulo organiza procedimentos estruturados de elicitação destinados a transformar saberes tácitos em variáveis linguísticas e critérios mensuráveis, assegurando consenso entre avaliadores e rastreabilidade plena das decisões. O método Delphi \cite{Linstone1975,Rowe1999} constitui técnica de elicitação de conhecimento estruturada, fundamentada em rodadas sucessivas de consulta a um painel de especialistas sob regime de feedback controlado e anonimato, visando convergência progressiva para consenso sobre variáveis caracterizadas por ambiguidade ou complexidade elevada. Em contextos de ativos tradicionais, o Delphi viabiliza a conversão de julgamentos qualitativos heterogêneos relativos à qualidade do manejo ou à autenticidade cultural em escalas operacionalizáveis e auditáveis.

\subsection{Protocolo de Implementação}

A seleção do painel envolverá o recrutamento de 15 a 25 participantes dotados de expertise diversificada em conhecimentos tradicionais, sistemas agroecológicos, gestão de propriedade intelectual e desenvolvimento territorial. A composição estimada contempla mestres de saberes quilombolas, pesquisadores acadêmicos, técnicos extensionistas e gestores públicos em proporções balanceadas. Os critérios de seleção abrangem experiência mínima de cinco anos na área de atuação, reconhecimento consolidado pela comunidade epistêmica ou territorial e disponibilidade para participar de três rodadas de consulta.

A primeira rodada, de caráter divergente, consistirá na aplicação de questionário aberto solicitando aos painelistas que enumeram variáveis consideradas relevantes para valoração de ativos tradicionais, organizadas segundo dimensões pré-estabelecidas englobando aspectos culturais, biofísicos, econômicos e simbólicos. As respostas serão consolidadas mediante análise de conteúdo sistemática, gerando lista preliminar de variáveis candidatas à incorporação no modelo.

Na segunda rodada, de natureza convergente, será enviado questionário estruturado contendo a lista consolidada de variáveis, solicitando aos painelistas que avaliem em escala Likert de cinco pontos a relevância de cada variável e proponham ajustes mediante fusões, exclusões ou adições. Para cada variável serão calculadas medianas, intervalos interquartis e coeficientes de variação, permitindo identificação de divergências remanescentes.

A terceira rodada, orientada ao consenso, envolverá reenvio do questionário acompanhado de feedback agregado contendo as medianas do grupo e o posicionamento individual de cada painelista na rodada anterior, solicitando reavaliação fundamentada. O consenso será operacionalmente definido como intervalo interquartil inferior ou igual a um ponto e mediana igual ou superior a quatro na escala de cinco pontos. Variáveis que não alcançarem consenso estatístico serão submetidas a análise qualitativa complementar para decisão fundamentada sobre manutenção ou exclusão.

A validação de conteúdo será conduzida mediante cálculo do Coeficiente de Validade de Conteúdo proposto por \citeonline{Hernandez-Nieto2002}, aplicado às variáveis que alcançaram consenso. Variáveis apresentando CVC igual ou superior a 0,80 serão consideradas validadas e incorporadas ao modelo fuzzy subsequente.

\subsection{Tratamento Analítico e Produtos}

Serão calculadas estatísticas descritivas englobando mediana, intervalo interquartil e moda para cada rodada, avaliando-se a evolução temporal do consenso mediante observação da redução progressiva do intervalo interquartil ao longo das rodadas sucessivas. A estabilidade das respostas será verificada mediante aplicação do coeficiente de concordância de Kendall, interpretando-se valores iguais ou superiores a 0,70 como indicativos de consenso forte.

Espera-se como produto um conjunto validado de variáveis linguísticas acompanhadas de definições operacionais acordadas pelo painel, as quais alimentarão subsequentemente a construção das funções de pertinência e regras de inferência do sistema fuzzy de valoração.
